\documentclass[paper=a4,fontsize=10.5pt,twoside=false,parskip=false]{scrartcl}
\usepackage[top=22mm,bottom=24mm,inner=22mm,textwidth=118mm,marginparwidth=42mm,marginparsep=8mm,headsep=10mm]{geometry}
\usepackage{fontspec}
\setmainfont{IBMPlexSerif}[Extension=.otf,UprightFont=*-Regular,ItalicFont=*-Italic,BoldFont=*-SemiBold,BoldItalicFont=*-SemiBoldItalic,Scale=0.94]
\setsansfont{IBMPlexSans}[Extension=.otf,UprightFont=*-Regular,ItalicFont=*-Italic,BoldFont=*-SemiBold,BoldItalicFont=*-SemiBoldItalic,Scale=0.94]
\newfontfamily\condensed{IBMPlexSansCondensed}[Extension=.otf,UprightFont=*-Regular,BoldFont=*-SemiBold,Scale=0.94]
\usepackage[english]{babel}\babelprovide[import,main]{romanian}
\usepackage{xcolor,microtype,booktabs,array,colortbl}
\usepackage[skins,breakable]{tcolorbox}
\definecolor{ink}{HTML}{1F2328}\definecolor{muted}{HTML}{6A737D}\definecolor{rule}{HTML}{D0D4D9}\definecolor{paper}{HTML}{FFFFFF}
\definecolor{accent}{HTML}{963A45}\definecolor{cex}{HTML}{4F7149}\definecolor{cnote}{HTML}{647080}\definecolor{cwarn}{HTML}{B0781F}\definecolor{csrc}{HTML}{74624F}
\pagestyle{empty}\setlength\parindent{1.2em}
\newcommand\Lbl[1]{{\sffamily\footnotesize\color{muted}\addfontfeatures{LetterSpace=8}\MakeUppercase{#1}}}
\input{\STYLE}
\begin{document}\color{ink}
\noindent\Lbl{\StyleName}\par\smallskip
\noindent{\small\color{muted}\StyleNote}\par\bigskip
\noindent Legea din 1864 a introdus învățământul primar obligatoriu și gratuit, dar aplicarea ei a depins de resursele fiecărei comune.
\begin{txbox}{accent}{Definiție}{Sursă primară}
Document produs în perioada studiată, de un participant sau martor al evenimentelor.
\end{txbox}
\begin{txbox}{cex}{Exemplu}{}
Raportul unui revizor școlar din 1872 privind școlile rurale din județul Vâlcea.
\end{txbox}
\begin{txbox}{cnote}{Notă}{}
Statisticile oficiale ale vremii supraestimează frecvent numărul elevilor înscriși.
\end{txbox}
\begin{txbox}{cwarn}{Atenție}{}
Nu confundați data adoptării legii cu data intrării ei în vigoare.
\end{txbox}
\begin{txbox}{accent}{Rezumat}{}
\begin{itemize}\setlength\itemsep{1pt}\item Învățământ primar obligatoriu din 1864. \item Rețeaua școlară se dublează până în 1914. \item Reformele Haret leagă școala de sat.\end{itemize}
\end{txbox}
\begin{txsource}{Legea instrucțiunii, 25 noiembrie 1864, art.~31}
„Instrucțiunea primară este obligatorie și gratuită. Părinții sau tutorii sunt datori a trimite copiii la școală de la vârsta de 8 până la 12 ani.”
\end{txsource}
\begin{txquote}{Nicolae Iorga, \emph{Istoria învățământului românesc}, 1928}
Școala satului nu poate fi mai bună decât învățătorul ei.
\end{txquote}
\end{document}
